\section*{Cluster A Step 36: Higher-Layer Refinement}

\paragraph{Grade.}
Finite-carrier diagnostic. The computation conjoins a staged scalar-potential compatibility predicate with the Step-35 higher-layer residual. It does not derive content, values, or a physical scalar sector.

\paragraph{Carrier.}
The Step-35 residual is rederived as \(12\) survivors.

\paragraph{Predicate.}
For each Step-35 scalar witness, the computation checks an invariant quadratic pairing, a bounded quartic proxy, and a staged residual-route condition: the scalar-active breaking route must leave at least one independent non-abelian route untouched.

\paragraph{Result.}
The predicate leaves \(8\) of \(12\) survivors. It eliminates the single-factor \(\texttt{4}\) family \((4 \to 0)\) and retains all \(8\) \(\texttt{2|3}\) survivors. The Step-14 reference support passes with witness \(\texttt{rank\_one\_fundxsinglet:-3}\), but it is not unique.

\paragraph{Verdict.}
\(\texttt{CONTENT\_TYPE\_LIMIT}\). The higher-layer descent grammar narrows to the target structure family in this finite carrier, but unique content selection remains outside this predicate.

\paragraph{Boundary.}
This is compatibility typing, not a derivation of content, scalar dynamics, empirical masses, or coupling values.
